%%%%%%%%%%%%%%%%%%%%%%%%%%%%%%%%%%%
% Cover Letter — Digital Discovery (RSC)
% Addressed to Prof. Alán Aspuru-Guzik, Editor-in-Chief
%%%%%%%%%%%%%%%%%%%%%%%%%%%%%%%%%%%

\documentclass[12pt,a4paper]{letter}
\usepackage[left=2.5cm, right=2.5cm, top=2.5cm, bottom=2.5cm]{geometry}
\usepackage{mathptmx}
\usepackage[T1]{fontenc}
\usepackage{microtype}
\usepackage{hyperref}
\usepackage[detect-all=true]{siunitx}

\signature{
  Jean-Pierre Tchapet Njafa\\
  (on behalf of all authors)
}

\address{
  Department of Physics, Faculty of Science\\
  University of Yaound\'{e}~1\\
  P.O. Box 812, Yaound\'{e}, Cameroon\\
  \texttt{jean-pierre.tchapet@facsciences-uy1.cm}
}

\date{June 27, 2026}

\begin{document}

\begin{letter}{
  Professor Alán Aspuru-Guzik\\
  Editor-in-Chief, \textit{Digital Discovery}\\
  Royal Society of Chemistry\\
  Thomas Graham House, Science Park\\
  Milton Road, Cambridge CB4 0WF, UK
}

\opening{Dear Professor Aspuru-Guzik,}

We are pleased to submit our manuscript entitled \textbf{``How Far Can
Structure-Based Machine Learning Predict Experimental Singlet--Triplet Gaps?
An Honest Benchmark for TADF Emitter Triage''} for consideration as a Research
Article in \textit{Digital Discovery}.

\textbf{Why \textit{Digital Discovery}?}
The journal has been a leading venue for rigorous, reproducible, and
self-critical machine learning in chemistry---including negative and
cautionary results. Our manuscript is in that spirit: a carefully validated
assessment of what cheap, structure-based ML can and cannot do for a property
of central interest, with every number reproducible from deposited code.

\textbf{Scientific contributions.}
The singlet--triplet gap ($\Delta E_\text{ST}$) governs thermally activated
delayed fluorescence (TADF) emitters. We make four contributions:

\begin{enumerate}
  \item \textbf{A diagnosis of feature--target leakage.} We show that regressing
  $\Delta E_\text{ST}$ onto feature sets containing the $S_1$ and $T_1$
  energies is circular ($\Delta E_\text{ST}\equiv E_{S_1}-E_{T_1}$); reported
  sub-\num{0.05}~eV accuracies for this set-up reflect arithmetic, not learning---a
  general leakage failure mode for any property defined as a difference of computed
  quantities.

  \item \textbf{A scaffold-validated benchmark.} On \num{231} molecules across
  \num{212} Bemis--Murcko scaffolds, a random forest on energy-free descriptors
  predicts \emph{experimental} $\Delta E_\text{ST}$ with MAE~$= 0.096$~eV
  ($R^2 = 0.25$, $\rho = 0.36$); a label-permutation test ($p = 0.001$) confirms the
  signal is real. Structural descriptors and fingerprints perform equivalently, and
  the semi-empirical excited-state step adds nothing.

  \item \textbf{An explicit accuracy ceiling.} Inter-laboratory scatter
  (\num{0.1}--\num{0.15}~eV) and functional dependence of the reference (up to
  \num{0.58}~eV between B3LYP and CAM-B3LYP) bound achievable accuracy near
  \num{0.1}~eV; a vertical-versus-adiabatic check (CAM-B3LYP/def2-TZVP) shows the
  cheap descriptors carry only a systematic, rank-preserving offset. The result has
  direct implications for how such models should be reported. As a robustness check, a
  threefold corpus expansion to \num{728} experimental emitters does not lower the error,
  confirming the ceiling is set by the measurement, not the sample size.

  \item \textbf{A practical triage filter} that enriches low-gap candidates by
  \num{1.2}--\num{1.4}$\times$ over random---useful for prioritisation, not a
  device-level or emission-spectrum predictor.
\end{enumerate}

\textbf{Novelty and differentiation.}
This work is distinct from our two prior \textit{JCIM} publications
(xTB/sTDA benchmarking and design-rule extraction): here we ask, rigorously,
what cheap structural data can and cannot support for $\Delta E_\text{ST}$
prediction, and we report the limitations openly.

\textbf{Open science.}
All computational data, trained models, SHAP analysis, and Python code are
openly available on Zenodo (DOI: 10.5281/zenodo.17436069) under CC-BY-4.0
license, ensuring full reproducibility and enabling extension to new
chemical spaces.

\textbf{APC waiver.}
As corresponding author affiliated with the University of Yaoundé I,
Cameroon, we are eligible for a 100\% APC waiver under the RSC's
Research4Life programme for authors from low- and middle-income countries.
We will provide the required documentation upon request.

\textbf{Suggested reviewers.}
We suggest the following experts who have the relevant background in ML
for materials discovery and TADF photophysics:
\begin{itemize}
  \item Prof. Heather Kulik (MIT) --- ML for molecular design, active learning
  \item Prof. Andrew Monkman (Durham) --- TADF photophysics, device performance
  \item Prof. Olexandr Isayev (Carnegie Mellon) --- ML for chemistry
  \item Prof. Denis Jacquemin (Nantes) --- TD-DFT for excited states
\end{itemize}

We confirm that this manuscript has not been published elsewhere and is not
under consideration by any other journal.
All authors have approved the manuscript and agree to its submission to
\textit{Digital Discovery}.

We hope you will find this work suitable for publication and look forward
to your response.

\closing{Yours sincerely,}

\end{letter}
\end{document}
